\subsection{A wild restriction and a sharp minimum in characteristic three}

\begin{lemma}[B-RD19: additive full odd torsion is potentially good]
\label{lem:B-RD19}
\status{Proved.}
Let $p$ be odd and let an elliptic curve $E/K$ have full rational
$p$-torsion. If its reduction is additive, it is potentially good.
Thus the additive existence question in B.3 is the potentially good
one, although its unrestricted good-forcing question also includes
multiplicative Tate curves.
\end{lemma}
\begin{proof}
The Weil pairing gives $\mu_p\subset K$. Suppose instead that
$v_K(j(E))<0$. The integral power series for $j(q)^{-1}$ starts with
$q$, so its compositional inverse gives $q\in K$, $v_K(q)>0$, with
$j(E_q)=j(E)$, where $E_q$ is the Tate curve
\cite[Theorem~1 and the $j$-expansion, pp.4--6]{Tate1995Uniformization}.
This value of $j$ is neither zero nor $1728$. Elliptic curves with
that same $j$ are quadratic twists, so $E=E_q\otimes\epsilon$ for
a quadratic character $\epsilon$. Tate uniformization gives
\[
       E[p]^{\mathrm{ss}}
        \simeq\overline\epsilon\,\overline\chi_p
                      \oplus\overline\epsilon.
\]
The left side is trivial and $\overline\chi_p=1$, so
$\overline\epsilon=1$. As $p$ is odd, reduction modulo $p$ is
injective on $\{1,-1\}$, hence $\epsilon=1$. The curve is therefore
split multiplicative, contrary to additive reduction. Thus $j(E)$
is integral, which is the elliptic criterion for potential good
reduction.
\end{proof}

\begin{lemma}[The maximum formal valuation is divisible by three]
\label{lem:B-RD20-wild}
\status{Proved; independent proof and source audits by the parent and obstructions/crystals agents passed.}
Suppose $p=3$, $E[3^n]\subset E(K)$, and $E$ has additive reduction
with wild semistability defect $m\in\{3,6,12\}$.
After extending the residue field algebraically, let $L/K$ be the
minimal good-reduction extension, with its normalized valuation
$v_L$. The maximum valuation $A$ of a primitive $3^n$-torsion
parameter in the good formal group is divisible by three.
\end{lemma}
\begin{proof}
First the discriminant of any Weierstrass equation for $E$ is a
cube in $K^\times$. The curve $E_0:y^2+y=x^3$ has discriminant
$-27=(-3)^3$ and full rational $3$-torsion over $K\supset\mu_3$,
as checked in Theorem~\ref{thm:B-RD2}. Choose a symplectic
isomorphism $E_0[3]\simeq E[3]$. The universal property of the
Hesse pencil makes $E$ $K$-isomorphic to a member of the Hesse
pencil of $E_0$ \cite[Theorem~1.1, p.2; proof, pp.7--8]
{AnemaTopTuijp2018Hesse}. Completing the square gives the short
equation $y^2=x^3+1/4$ with the same discriminant. Proposition~5.1
of that source, pp.4--5, computes the discriminant of every finite
smooth member as $\Delta(E_0)(\det A_t)^3$; the member at infinity
is $E_0$ itself. Thus every such discriminant is a cube.
A change of Weierstrass coordinates changes the discriminant by a
twelfth power. Hence $\delta=v_K(\Delta)$ is divisible by three.

Write the relation between a $K$-rational invariant differential and
an integral generator on the good model over $L$ as
$\omega_L=u\omega_K$. Since the good discriminant is a unit,
\[
                  v_L(u)=m\delta/12.
\]
This is an integer divisible by three, for each
$m=3,6,12$. Let $\sigma$ generate the wild subgroup $C_3$, and set
$F=L^{C_3}$. Choose $b_0\in F^\times$ with
$v_L(b_0)=v_L(u)$. Replacing the integral differential by
$(b_0/u)\omega_L$ makes it $C_3$-invariant without changing its
integrality, since $b_0/u$ is a unit. Choose a formal parameter with
this linear differential and odd under $[-1]$, which is possible
because two is a unit. Thus the wild descent series has linear
coefficient exactly one and only odd powers.

The special fibre is supersingular by Proposition~\ref{prop:B-RD3}.
On $y^2=x^3-x$ its order-three automorphisms are translations
$x\mapsto x+r$, $r\in\mathbf F_3^\times$. In the odd parameter
$t=-x/y$, their first nonidentity term has degree three and nonzero
coefficient. Therefore the wild descent series, or its inverse
according to the coordinate convention, has the form
\[
                   t\longmapsto t+c t^3+O(t^5),
                   \qquad c\in\mathcal O_L^\times.
\]
For a $K$-rational nonzero torsion point with parameter $x$ this gives
\[
                     v_L(\sigma x-x)=3v_L(x).       \tag{B-RD20.1}
\]

Let $b=v_L(\sigma\pi-\pi)-1$ be the positive lower ramification
break, for a uniformizer $\pi$ of $L$. The standard uniformizer
expansion, with coefficients in the fixed Cohen ring $W(k)$, gives
\[
 v_L(\sigma x-x)\geq v_L(x)+b,
 \quad\text{with equality if }3\nmid v_L(x).       \tag{B-RD20.2}
\]
Indeed the term of lowest valuation in $x$ gives equality when its
exponent is prime to three, and every higher term has strictly
larger valuation in its difference. If $B$ is the minimum primitive
valuation, (B-RD20.1)--(B-RD20.2) imply $b\leq2B$. If $A>B$ and
$3\nmid A$, the same equations give $b=2A>2B$, a contradiction.

It remains to exclude $A=B$ with $3\nmid A$. In that case $b=2A$.
For every primitive torsion parameter $x$, write
$x=a_x\pi^A+O(\pi^{A+1})$, with $a_x\in k^\times$. The first
coefficient of $\sigma x-x$, divided by $\pi^{3A}$, is a fixed
nonzero scalar times $a_x$; the first coefficient of the right
side of the descent equation is $\bar c a_x^3$. Thus all these
coefficients solve one equation
\[
                           \bar c X^3=dX,
                           \qquad d\ne0.
\]
Here $d=A\overline{(\sigma\pi-\pi)/\pi^{b+1}}\ne0$, since
$3\nmid A$. Its roots have dimension at most one over
$\mathbf F_3$. Choose a
primitive basis $P,Q$ of $E[3^n]$. Sending $(a,b)\in\mathbf F_3^2$
to the leading coefficient of $[\widetilde a]P+[\widetilde b]Q$
is well defined and linear: changing an integer lift by a multiple
of three does not affect the leading coefficient, because formal
addition is addition to first order and the residue characteristic
is three. This map is injective, since a nonzero combination is
primitive and cancellation would give valuation greater than the
maximum $A$. This requires dimension two, a contradiction. Hence
$3\mid A$ also in this case.
\end{proof}

\begin{proposition}[B-RD20: a stricter wild bound and an exact minimum]
\label{prop:B-RD20}
\status{Proved; independent proof and source audits by the parent and obstructions/crystals agents passed.}
For $p=3$ write $e=2u$ and $s=v_3(u)$. If full rational $3^n$-torsion
occurs with wild defect $m\in\{3,6,12\}$, necessarily
\[
 s\geq n-1,\qquad
 e\geq\frac8m\,3^{\max(2n-2,\,4n-4-s)}.        \tag{B-RD20.3}
\]
Allowing every perfect residue field and every defect, the least
absolute ramification index admitting an additive elliptic curve
with full rational $3^n$-torsion is exactly
\[
                            e_{\min}(n)=2\cdot3^{2n-2}.
                                                        \tag{B-RD20.4}
\]
In particular $e_{\min}(2)=18$: absolute indices six and twelve,
although compatible with the cyclotomic necessary condition,
do not admit an additive curve with full rational $9$-torsion.
\end{proposition}
\begin{proof}
The cyclotomic necessary condition gives $s\geq n-1$. The
valuation identity in the proof of Theorem~\ref{thm:B-RD12} uses
only the good height-two formal group, not the tame descent
congruences. In a valuation region
$3^{2r}B\leq A\leq3^{2r+2}B$, $0\leq r\leq n-1$, it therefore
still gives
\[
                  me=2\cdot3^{2n-r-2}(A+3^{2r+1}B).
                                                        \tag{B-RD20.5}
\]
The upper bound for $A$ is omitted when $r=n-1$.
By Lemma~\ref{lem:B-RD20-wild}, the parenthesis is divisible by
three. Since $v_3(m)=1$, this forces
$s\geq2n-r-2$, or $r\geq\max(0,2n-2-s)$.
The lower bound on $A$ and $B\geq1$ give
$me\geq8\cdot3^{2n+r-2}$. Substitute the least possible $r$ to
obtain (B-RD20.3).

To prove (B-RD20.4), tame defects satisfy the lower bound
$e\geq2\cdot3^{2n-2}$ by Theorem~\ref{thm:B-RD12} with $m\leq4$.
For a wild defect use $m\leq12$ in (B-RD20.3). If
$s\geq2n-2$, divisibility of $e$ by $2\cdot3^s$ already gives
the claimed bound. Otherwise $s\leq2n-3$, and the second exponent
in (B-RD20.3) is at least $2n-1$. Consequently
\[
 e\geq\frac23\,3^{2n-1}=2\cdot3^{2n-2}.
\]
Finally the tame construction with $m=4$ and
$e=2\cdot3^{2n-2}$ meets the criterion of
Theorem~\ref{thm:B-RD12}, so it attains this bound for every $n$.
Lemma~\ref{lem:B-RD19} covers every additive curve in the question.
\end{proof}

\begin{corollary}[B-RD21: the remaining wild question is one level]
\label{cor:B-RD21}
\status{Proved; independently audited by the parent agent; arithmetic certificate supplied.}
Fix $p=3$ and a positive even absolute ramification index $e$.
Let $N(e)$ be the largest $n\geq1$ for which an additive elliptic
curve with full rational $3^n$-torsion exists over some such field
with perfect residue field. Put
\[
 s=v_3(e/2),\qquad L=\lfloor\log_3(e/2)\rfloor,
\]
and define
\begin{align*}
 N_{\mathrm t}
   &=1+\min\left(s,\left\lfloor\frac L2\right\rfloor,
                         \left\lfloor\frac{L+s}{4}\right\rfloor\right),\\
 N_{\mathrm w}
   &=1+\min\left(s,\left\lfloor\frac{L+1}{2}\right\rfloor,
                         \left\lfloor\frac{L+s+1}{4}\right\rfloor\right).
\end{align*}
Then
\[
             N_{\mathrm t}\leq N(e)\leq N_{\mathrm w}
                             \leq N_{\mathrm t}+1.       \tag{B-RD21.1}
\]
In particular the full answer is already determined whenever
$N_{\mathrm t}=N_{\mathrm w}$. A gap occurs precisely when there is
an integer $k\geq1$ such that
\[
             k\leq s\leq2k-1,\qquad L=4k-s-1.          \tag{B-RD21.2}
\]
Then $N_{\mathrm t}=k$ and $N_{\mathrm w}=k+1$, so only existence
at level $n=k+1$ remains undecided. For full $9$-torsion, the only
indices not decided by these results are
\[
                              e=24,30,42,48.            \tag{B-RD21.3}
\]
At $e=24,30$ an example, if one exists, must have defect twelve;
at $e=42,48$ its possible wild defects are six and twelve.
\end{corollary}
\begin{proof}
The sharp tame criterion with defect four gives the lower bound
$N_{\mathrm t}$. The bound in Proposition~\ref{prop:B-RD20} is
weakest for defect twelve; its logarithmic inequalities are
$n-1\leq s$, $2n-2\leq L+1$, and $4n-4-s\leq L+1$.
They give the upper bound $N_{\mathrm w}$. Increasing $L$ by one
increases each displayed minimum by at most one.

Since $s\leq L$, a level $k+1$ admitted by the wild inequalities
but excluded by the tame ones must satisfy
\[
 k\leq s,\qquad \max(2k,4k-s)=L+1.
\]
The maximum cannot be $2k$ with $s\geq2k$, because that would
give $L=2k-1<s$. Thus $s\leq2k-1$ and $L=4k-s-1$.
Conversely these conditions make the two minima respectively
$k-1$ and $k$, which proves (B-RD21.2).
For $k=1$ we have $s=1,L=2$, so $e/2$ is a multiple of three
but not nine in the interval $[9,27)$; the four possibilities
are $12,15,21,24$. Substitute them into (B-RD20.3) for the
defect assertions. The finite integer certificate
\path{code/verify_p3_higher_wild_bounds.py} checks these formulae,
their extrema, and their compatibility with the unsimplified
valuation identity in the stated finite range.
\end{proof}
